\section{模型假设}

为在保证物理合理性的前提下使问题可解，本文做出如下基本假设。凡涉及题目未给出显式公式之处，均在相应小节标注"建模口径"并说明依据，不擅自代入未声明的规则。

\begin{enumerate}
    \item \textbf{水平巡航能耗由等效航程包络反推}。运输机规格表仅给出空载/满载标准航程与电池可用能量，未给出巡航功率列，故取单位距离能耗 $E^{\rm hor}(d,q)=d/L_g(q)\cdot E^{\rm use}_g$，即"满电可飞 $L_g(q)$ 距离"反推单位距离能耗；爬升附加能耗按重力势能除以爬升效率计算，下降附加能耗按题设取 0。
    \item \textbf{航段按爬升—巡航—下降三段计算}。计划巡航海拔取航段所经 DEM 像元最高地面高程以上 50\,m；O01 作业高度取地面海拔，服务区作业高度取地面海拔以上 30\,m；爬升/下降段视为在起讫点水平位置原地垂直升降，巡航段为定高直线飞行。
    \item \textbf{货箱不可拆分且每箱仅运输一次}。组批时同服务区内的货箱可任意组合，但不得跨服务区组批；单箱质量/体积不超过任一机型上限（问题一已验证单箱必可被某机型单独运输）。
    \item \textbf{共享电池与机身均为可互换资源池}。同机型电池可在该机型不同实体机间调度、跨机型不可混用；无人机架次结束后机身立即就绪，电池结束后进入两阶段非线性充电；电池初始 SOC 均为 100\%，再次投入前须充至 100\%。
    \item \textbf{多站架次停靠服务区数上限为 3}。问题一 $N=18$ 架次/15 区，多数服务区单独接近满载，3 站以上同时满足载重的概率极低，属有依据的搜索空间裁剪。
    \item \textbf{通信链路按自由空间损耗 + 二元地形遮挡判定}。路径损耗采用 ITU-R P.525 自由空间公式\cite{ref_itu}，地形遮挡按 DEM 视线二元判定并加遮挡附加损耗 $L_{\rm obs}$，不做渐进式菲涅尔区计算；链路按双向（取两个方向门限较小者）判定。
    \item \textbf{中继仅支持单跳}。运输无人机经一架中继建立"运输—中继—G01"两段链路，不允许中继间多跳转发（数据中无中继—中继参数，属数据结构决定）。
    \item \textbf{中继悬停位置离散化}。中继悬停位置在"缺口腿水平连线的若干分数点 $\times$ 多档悬停高度"候选网格上搜索，非连续优化；每次中继任务视为"O01 起飞—到位服务—飞回 O01"的完整架次。
    \item \textbf{相位内任一采样点被遮挡即判定该相位全程需中继覆盖}。此判定偏保守，但避免了子相位区间的复杂追踪，且能保证通信连续性约束不遗漏。
    \item \textbf{题目所给数据真实可靠}。空间坐标与公开地理数据仅用于描述测试场景，需求、时限、装备与资源参数以附件给定值为准；计算全程保留原始精度不截断。
\end{enumerate}

上述假设中，假设 1--2 统一了全文能耗与时间口径，假设 4、5 定义了资源调度与搜索边界，假设 6--9 界定了通信与中继的判定口径，假设 10 保证数据口径一致。各假设的影响范围及可能带来的偏差将在对应小节与模型评价中说明。
